% Generated draft: requires mathematical and source review.
\documentclass[11pt]{article}
\usepackage[margin=25mm]{geometry}
\usepackage{fontspec}
\setmainfont{Latin Modern Roman}
\usepackage{amsmath,amssymb,mathtools,unicode-math}
\setmathfont{Latin Modern Math}
\usepackage{enumitem,needspace,xurl,hyperref}
\setlength{\parindent}{0pt}
\setlength{\parskip}{5pt}
\setlength{\emergencystretch}{3em}
\newcommand{\N}{\mathbb{N}}
\newcommand{\Z}{\mathbb{Z}}
\newcommand{\Q}{\mathbb{Q}}
\newcommand{\R}{\mathbb{R}}
\newcommand{\C}{\mathbb{C}}
\newcommand{\F}{\mathbb{F}}
\newcommand{\A}{\mathbb{A}}
\newcommand{\PP}{\mathbb{P}}
\newcommand{\E}{\mathbb{E}}
\newcommand{\eps}{\varepsilon}
\newcommand{\dd}{\,\mathrm{d}}
\begin{document}
\section*{UnsolvedMath ID: 701705}
\textbf{Source problem code:} catalogue entry
\subsection{Definitions and mathematical statement}
Let \(G=F_n\) be a free group of finite rank, regarded as a discrete group. The supplied statement imposes no restriction \(n\geq 2\). Under the usual convention for a compact \(G\)-space, let \(X\) be a nonempty compact Hausdorff space with a continuous action \(G\times X\to X\). These topological conventions are implicit rather than explicit in the supplied record.

The action is \emph{faithful} if
\[
\{g\in G: gx=x\text{ for every }x\in X\}=\{1_G\},
\]
and \emph{minimal} if \(\overline{Gx}=X\) for every \(x\in X\). Let \(\mathcal P(X)\) denote the regular Borel probability measures on \(X\), with the weak-* topology, and define the pushforward action by \((g_*\mu)(A)=\mu(g^{-1}A)\) for Borel sets \(A\subseteq X\). The action on \(X\) is \emph{strongly proximal} if
\[
\forall\mu\in\mathcal P(X)\quad\exists x\in X\quad
\delta_x\in\overline{\{g_*\mu:g\in G\}}^{\,\mathcal P(X)},
\]
where \(\delta_x\) is the point mass at \(x\). It is \emph{topologically free} if
\[
\forall g\in G\setminus\{1_G\},\qquad
\operatorname{int}_X\{x\in X:gx=x\}=\varnothing.
\]
\emph{Question (Kourovka Notebook Problem 20.47).} For every finitely generated free group \(G=F_n\) and every compact \(G\)-space \(X\), under the conventions above, does faithfulness, minimality, and strong proximality imply that the action is topologically free?
\subsection{Short English statement}
Must every faithful, minimal, strongly proximal action of a finitely generated free group on a nonempty compact Hausdorff space have the property that no nonidentity element fixes a nonempty open subset? The compact Hausdorff and nonemptiness conventions are implicit in the supplied source record.
\subsection{Sources}
\begin{itemize}
\item Local UnsolvedMath registry: permanent ID \(701705\), exact source problem code \texttt{catalogue entry}, titled ``Kourovka Notebook Problem 20.47.'' The supplied record labels the problem open; this label is supplied curation, not an independently verified mathematical or literature assessment.
\item Attributed to A. Le Boudec and N. Matte Bon, proposed in 2022. The supplied background identifies \emph{The Kourovka Notebook (Unsolved Problems in Group Theory)}, 20th issue (2022), problem 20.47, and the 21st issue edition (2026), arXiv:1401.0300v48: \url{https://arxiv.org/abs/1401.0300v48}. The supplied collection description names E. I. Khukhro and V. D. Mazurov as editors.
\item The supplied registry record was updated on \texttt{2026-10-16T14:00:00Z} and identifies upstream revision \texttt{d371ee917ac9647fbb06cd6fc529deec6f3dda82}. This timestamp is registry metadata, not a supplied literature-verification date. The cited source has not been independently consulted for this draft.
\end{itemize}
\end{document}
